\section{Problem setup}
\label{sec:setup}

\subsection{Units, risk and the guarantee}

Let $Z=(X,A^\star)$ be a query--answer pair, where $X$ contains an image and a question and
$A^\star\in\{0,1\}$ is the gold answer. A VLM returns $\hat a(X)$. A second predictor (channel 2)
returns an answer $b(X)$ and an evidence margin $m(X)\ge0$. Write $Y=\ind{\hat a(X)=A^\star}$ and
$Y^{(2)}=\ind{b(X)=A^\star}$ for the two correctness indicators and $\mu=\Pr[Y=0]$ for the base
risk. $\alpha$ is always a risk target and $\delta$ the failure probability of the certificate.

A \emph{policy} $\pi$ maps $X$ to one of three actions: emit $\hat a(X)$, emit $b(X)$, or abstain.
Let $\emit_\pi(X)=\ind{\pi\text{ emits}}$ be the emission indicator and
$W_\pi=\emit_\pi\cdot\ind{\text{emitted answer}\ne A^\star}$ the indicator of an emitted error. The
\emph{risk} and \emph{coverage} of $\pi$ are
\[
\Risk(\pi)=\frac{\E[W_\pi]}{\Pr[\emit_\pi=1]},\qquad \Cov(\pi)=\Pr[\emit_\pi=1],
\]
with the convention $\Risk(\pi)=0$ when $\Pr[\emit_\pi=1]=0$. The same convention applies to a
realised abort: a procedure that certifies nothing emits nothing, has risk $0$ and coverage $0$,
and every coverage figure in this paper includes aborted splits as zero coverage.

\begin{assumption}[Sampling unit]
\label{ass:iid}
The calibration units $Z_1,\dots,Z_n$ are i.i.d.\ draws from a distribution $P$, independent of the
data used to build the policy family (the \emph{fit} data), and the unit is the item.
\end{assumption}

Exchangeability is not enough here. Under exchangeability alone, sampling without replacement from
a finite pool makes the number of emitted errors hypergeometric rather than binomial, and the exact
binomial bound used below is not valid. When several questions share an image the item is not an
i.i.d.\ unit; \S\ref{sec:clusters} states what remains true and measures what does not.

\begin{assumption}[Pre-specified policy family]
\label{ass:family}
The candidate policies form a finite, ordered family $\pi_1,\dots,\pi_M$ whose members and order
are functions of the fit data and of constants only, never of calibration outcomes or calibration
features.
\end{assumption}

Calibration is a statement about the \emph{population} risk $\Risk(\pi_{\hat\jmath})$ of the policy
the calibration sample selects; it is not a statement about any finite test fold, whose realised
error rate carries binomial noise of its own.

\subsection{Filtering and the floor it cannot cross}

A \emph{filtering} policy uses a score $s(X)\in[0,1]$, large when the model is likely right, and
emits $\hat a(X)$ when $s(X)\ge t$.

\begin{proposition}[Abstention floor; Kotte \citep{crcimposs}]
\label{prop:floor}
Let $\mu>\alpha$. Any selective predictor whose emitted answer is always $\hat a(X)$ and whose risk
is at most $\alpha$ satisfies
\[
1-\Cov\ \ge\ \frac{\mu-\alpha}{1-\alpha}.
\]
\end{proposition}
\begin{proof}
Split the error mass by emission, $\mu=\Risk\,\Cov+\Pr[Y=0\mid\text{abstain}](1-\Cov)\le\alpha\Cov+(1-\Cov)$.
The inequality uses that abstaining does not change $Y$. Solving for $\Cov$ gives
$\Cov\le(1-\mu)/(1-\alpha)$.
\end{proof}

The proof uses one premise: an abstained item's model answer is still wrong. It is a statement about
the action space. Because $\alpha$ is chosen by the user and $\mu$ is estimated from data, the
floor can be evaluated before calibration only up to the uncertainty of $\hat\mu$; a feasibility
check should use a confidence bound on $\mu$, not its plug-in value.

\subsection{Two channels}

\begin{definition}[Channels]
\label{def:channels}
\emph{Channel 1} is a correctness score $s(X)\approx\Pr[Y=1\mid\phi(X)]$ fitted on features of the
model's own output together with grounding features. \emph{Channel 2} is a distinct predictor
$b(X)$ that answers the same question without using the VLM's answer, together with a margin
$m(X)$ measuring how decisively it does so.
\end{definition}

We call channel 2 a \emph{distinct evidence source}. Statistical independence between $Y$ and
$Y^{(2)}$ is neither assumed nor needed for validity: Theorem~\ref{thm:validity} holds for any
fixed policy and needs only the sampling model. What distinctness buys is usefulness. A repair
region helps only if $\Pr[Y^{(2)}=1\mid\text{region}]$ is high, which requires channel 2's errors to
differ from the model's inside that region (\S\ref{sec:selfrepair} shows what happens when channel 2
is the negation of the model's answer). Our instantiation also violates independence in the
other direction: the detector score is an input to channel 1.

\paragraph{Instantiation}
Channel 2 is OWLv2 \citep{owlv2} queried with the object named in the question. With detection
confidence $o(X)$ and a fixed operating point $\theta=0.15$, $b(X)=\ind{o(X)\ge\theta}$ and
$m(X)=|o(X)-\theta|$. A confident negative reaches only $m\le\theta$ while positives reach
$1-\theta$, so with a single gate on $m$ every repair says ``the object is present''. Section
\ref{sec:accounting} evaluates a polarity-symmetric gate as well. When no detector value exists
(the object phrase could not be parsed, or the question concerns an attribute or relation), the
item keeps channel 1's decision, receives $m=-1$, and can never be repaired.

\begin{table}[t]
\centering
\caption{Notation.}
\label{tab:notation}
\small
\begin{tabular}{@{}lp{0.78\textwidth}@{}}
\toprule
$\alpha,\delta$ & risk target; failure probability of the certificate\\
$\mu$ & base risk $\Pr[Y=0]$\\
$Y,\ Y^{(2)}$ & correctness indicators of the model and of channel 2\\
$s(X),\ m(X)$ & channel-1 score; channel-2 evidence margin\\
$t_1>\dots>t_M$ & acceptance thresholds, score \emph{values} fixed from fit data\\
$g$ & repair-gate value on $m$, fixed from fit or external data\\
$\emit_j$, $K_j$, $V_j$ & emission indicator, emitted count and emitted-error count of policy $j$ on the calibration sample\\
$U(e,k;\delta)$ & one-sided Clopper--Pearson upper bound\\
$k_{\mathrm{start}}(e_0)$ & smallest $k$ with $U(e_0,k;\delta)\le\alpha$\\
\bottomrule
\end{tabular}
\end{table}
